\chapter{The Executable Certificate}
\label{ch11:chapter}

\section{Why an executable layer exists}
\label{ch11:sec:why}

The finite-context theory of the preceding chapters has an unusual property
for a body of category-theoretic results: every quantifier in every statement
ranges over finite, small, explicitly enumerable data. There are $4$ atoms,
$16$ contexts, $81$ inclusion arrows, and --- as we derive below --- exactly
$298$ sieves in total across all contexts. The witness presheaves have at most
four sections at any context. Nothing is infinite, nothing is chosen, nothing
depends on a universe beyond bookkeeping. A theory this small can be checked
\emph{exhaustively}: not sampled, not spot-checked, but walked in its
entirety by a machine.

The Noetic Topos 2.0 candidate does exactly that, with two independent
executable layers and a bounded proof-assistant layer. The first layer is a
dependency-free finite-site checker whose output is the registered artifact
\texttt{FINITE\_MODEL\_CERTIFICATE\_v0.1.json}, covering the site
combinatorics, subcanonicity, the witnesses, the marked-cover integrity
rules, and both legacy collapse mechanisms. The second is the
modality-composition verifier, which implements the explicit reflector of
Chapter~10 and computes both relevant Grothendieck topologies \emph{exactly},
by enumeration and saturation, rather than trusting the pen-and-paper join
argument. The third is a deliberately narrow Lean scope. This chapter walks
through all three, check by check, and closes with the certificate's own
statement of what all this execution does and does not establish.

\begin{boundary}
\label{ch11:bnd:role}
The relationship between the proofs and the programs must be stated with
care, because it is easy to overclaim in either direction. The chapters'
informal proofs are general arguments over the fixed finite site; the
programs are exhaustive executions over the same site. Execution
\emph{corroborates} the proofs --- a bug in a proof about a $16$-object site
would almost certainly produce a counterexample the exhaustive run would
find --- but it is not itself a proof: the checker could share a
misformalization with the prose, and passing tests are evidence about the
implemented model, not about the mathematical English. Conversely the proofs
do not make the programs redundant: the programs pin down every convention
(indexing, empty pieces, null covers, serialization) at a level of
concreteness prose cannot reach, and they are what a future reviewer can
rerun. Neither layer is licensed to substitute for the other, and neither is
licensed to substitute for independent review.
\end{boundary}

\section{Methodology of exhaustive checking}
\label{ch11:sec:methodology}

Before walking the checks, four methodological commitments deserve
statement, because they distinguish this executable layer from ordinary
software testing and because each is auditable in its own right.

\emph{Determinism.} Both verifiers are deterministic programs over fixed
finite data. There is no randomness, no sampling, no property-based
generation with a seed: a run either passes identically every time at a
given commit or fails identically every time. This is what allows the
certificate to be a \emph{record} rather than a statistic --- the
expected-result block asserts exact values, not distributions, and a
reviewer who reruns the registered command either reproduces the block
bit-for-bit or has found a discrepancy worth filing.

\emph{Dependency-freedom.} The finite-site checker is deliberately free
of third-party dependencies: contexts are frozen sets, presheaves are
dictionaries, sieves are sets of subsets, and every categorical notion is
implemented from its definition rather than imported from a library. The
cost is verbosity; the benefit is that the trusted computing base for the
certificate is the checker file itself plus the language runtime, and a
reviewer auditing the formalization has one file to read against
Chapter~10's definitions.

\emph{Exhaustive quantification.} Every ``for all'' in the checked
properties is discharged by enumeration over its entire domain: all
contexts, all inclusion arrows, all generating covers, all members, all
sieves where sieves are quantified. Section~\ref{ch11:sec:domains}
tabulates the domain of each check explicitly, because ``exhaustive'' is
only meaningful relative to a stated domain.

\emph{Negative expectations as first-class results.} The certificate
registers expected \emph{failures} --- the two false sheaf verdicts, the
rejection branches of the resolver, the two legacy collapse mechanisms
--- with the same standing as expected successes. A model in which
$F_A$ accidentally satisfied RPM gluing would \emph{fail} the
certificate, exactly as a model in which a representable failed the
sheaf condition would. This matters for evidential value: a check suite
that can only confirm positives cannot catch a construction that is
degenerately too strong, and the historical failure being repaired ---
the legacy collapse --- was precisely a construction whose covers were
degenerately too strong (maximal everywhere).

\section{The finite-site checker and its twenty-three checks}
\label{ch11:sec:finitesite}

The finite-site layer is the module \texttt{src/noetic\_topos2/finite\_site.py}
with test suite \texttt{tests/noetic\_topos2/test\_finite\_site.py}. Its
certificate declares schema version
\texttt{noetic\_topos2.finite\_model\_certificate.v0.1}, document status
\textsc{internal finite check complete / independent proof review open}, and
model identifier
\begin{center}
\texttt{NT2-CONTEXT-2X2-TRANSVERSE-PARTITIONS},
\end{center}
naming exactly the model of the monograph: the $2 \times 2$ atom grid, its
powerset context lattice, and the two transverse coordinate partitions. The
certificate's expected-result block registers two cardinality checks and
twenty-one named verdicts (counting each paired ACBE/RPM verdict as a single
named check), twenty-three registered checks in all, together with an
aggregate flag \texttt{all\_required\_checks\_pass}. All pass at the bound
commit. We take them in order, stating for each what is computed and which
result of the monograph it corroborates.

\subsection{Cardinalities and cover geometry}
\label{ch11:sec:geometry}

\begin{checkclass}
\label{ch11:chk:cardinalities}
\textbf{\texttt{atom\_count} $= 4$, \texttt{context\_count} $= 16$.} The
checker constructs $\Omega$ and $L = \mathcal{P}(\Omega)$ from scratch and
counts. These two numbers anchor every later enumeration: a model that
silently dropped an atom or a context would fail here before any
sheaf-theoretic check ran. Corroborates the standing data of the context
lattice definition (the sixteen-object poset with $81$ inclusion arrows; the
arrow count $\sum_{V} 2^{|V|} = 1 + 4 \cdot 2 + 6 \cdot 4 + 4 \cdot 8 + 16 =
81$ is exercised by the composition verifier below).
\end{checkclass}

\begin{checkclass}
\label{ch11:chk:proper}
\textbf{\texttt{proper\_global\_covers} (ACBE and RPM).} Verifies that at the
global context $\Omega$, every member of $K_A(\Omega) = \{A_0 \to \Omega,
A_1 \to \Omega\}$ and of $K_R(\Omega)$ is a \emph{proper} subcontext. This is
the single most load-bearing geometric fact in the repair: the legacy
clauses collapsed precisely because their generating sieves contained
identities, and a sieve on a poset generated by arrows from proper subobjects
cannot contain the identity. Corroborates the properness proposition of the
finite-context theory (ledger \textsc{thm-nt2-nontrivial-global-covers}) and,
by contrast, the collapse diagnosis for the legacy site.
\end{checkclass}

\begin{checkclass}
\label{ch11:chk:exhaustive}
\textbf{\texttt{jointly\_exhaustive} (ACBE and RPM).} Verifies
$A_0 \cup A_1 = \Omega = R_0 \cup R_1$ and, contextwise, that the pieces of
$K_P(V)$ union to $V$. Joint exhaustiveness is what makes the generating
families \emph{covers} in the intuitive sense and is used in the
subcanonicity proof (a compatible family over all pieces determines behavior
on all of $V$). Corroborates the coverage-basis lemma behind Theorem NT2-1.
\end{checkclass}

\begin{checkclass}
\label{ch11:chk:pullback}
\textbf{\texttt{pullback\_stable} (ACBE and RPM).} For every pair
$U \subseteq V$ and both partitions, the checker pulls back every member of
$K_P(V)$ along the inclusion and verifies the resulting family is $K_P(U)$
(as an indexed family; identity arrows on the empty object when
$U = \varnothing$, where the registered null cover covers as well). Pullback
stability is axiom (J2) at the level of the basis and is the hypothesis under
which basis-level gluing suffices for the generated topology
\cite{stacks00ZC}. Corroborates the coverage-basis lemma and hence Theorem
NT2-1 (ledger \textsc{thm-nt2-marked-coverage}).
\end{checkclass}

\subsection{Subcanonicity}
\label{ch11:sec:subcanonicity}

\begin{checkclass}
\label{ch11:chk:representables}
\textbf{\texttt{representables\_are\_sheaves}.} For each of the $16$ objects
$W \in L$, the checker builds the representable presheaf $yW$ (a singleton
at $V$ when $V \subseteq W$, empty otherwise), and verifies the gluing
bijection at every context against every generating cover of both
topologies, including the null cover at the empty context. This is the full
finite content of Theorem NT2-2 (ledger
\textsc{thm-nt2-subcanonical-context-factor}): both $J_A$ and $J_R$ are
subcanonical on $L$. It is worth noting what a failure here would mean: a
representable failing the sheaf condition would put the finite site in the
same position as the legacy site, whose jointly-surjective coverage is
provably not derivably subcanonical; the check certifies that the repair
does not reintroduce that gap on the context factor.
\end{checkclass}

\subsection{Witness functoriality and the four verdicts}
\label{ch11:sec:witnesses}

\begin{checkclass}
\label{ch11:chk:functorial}
\textbf{\texttt{acbe\_witness\_functorial}, \texttt{rpm\_witness\_functorial}.}
Before any sheaf verdict is computed, the checker validates that the proposed
witness data $F_A$ and $F_R$ --- bit assignments on complete contained
blocks, with restriction forgetting coordinates --- actually constitute
presheaves: identity restrictions are identities and all composable
restriction pairs commute. This ordering is deliberate methodology:
a nonfunctorial section assignment could otherwise pass or fail sheaf
conditions meaninglessly. The check corroborates the well-definedness half of
the distinctness theorem's witness construction (Theorem NT2-3), and it
enforces the program-wide discipline that a ``witness'' must be validated as
an object before being scored as evidence.
\end{checkclass}

\begin{checkclass}
\label{ch11:chk:verdicts}
\textbf{The four sheaf/non-sheaf verdicts:
\texttt{acbe\_witness\_is\_acbe\_sheaf} (true),
\texttt{acbe\_witness\_is\_rpm\_sheaf} (false),
\texttt{rpm\_witness\_is\_rpm\_sheaf} (true),
\texttt{rpm\_witness\_is\_acbe\_sheaf} (false).} The checker computes, for
each witness and each topology, the matching-family set of every generating
cover at every context and compares it with the value set. The registered
expectations are asymmetric on purpose: two verdicts are expected
\emph{true} and two expected \emph{false}, and the certificate records the
false ones as first-class expectations rather than as absences. Together the
four verdicts are exactly the content of Theorem NT2-3 (ledger
\textsc{thm-nt2-distinct-modalities}) as scoped to $\PSh(L, \Set)$: each
subcategory $\Sh(L, J_A)$, $\Sh(L, J_R)$ contains an explicit object the
other excludes, so the subcategories are distinct and neither modality is
trivial. The quantitative core --- $F_A(\Omega)$ has $4$ sections while the
RPM matching-family set at $\Omega$ is a singleton, so the gluing map
$4 \to 1$ fails injectivity --- is tabulated in full in Appendix~C.
\end{checkclass}

\subsection{Transverse refinement and singleton contexts}
\label{ch11:sec:refinement}

\begin{checkclass}
\label{ch11:chk:transverse}
\textbf{\texttt{transverse\_refinement\_is\_singletons}.} Computes all
sixteen intersections $A_i \cap R_j$ over the four index pairs and verifies
each equals the singleton $\{(i, j)\}$; contextwise, verifies that the
common refinement of $K_A(V)$ and $K_R(V)$ has as nonempty members exactly
the singletons of $V$. This is the transversality identity, the one equation
on which the join lemma, the composition theorem's flattening bijection, and
the four-channel expressivity boundary all rest. Corroborates Theorem NT2-4
(ledger \textsc{thm-nt2-join-discrete-context}) and the input identity of
Chapter~10's Theorem on modality composition.
\end{checkclass}

\begin{checkclass}
\label{ch11:chk:singleton}
\textbf{\texttt{singleton\_unmarked\_families\_coincide},
\texttt{singleton\_marked\_covers\_distinct}.} At each singleton context
$\{(i, j)\}$, the checker verifies two facts in deliberate tension. First,
the bare arrow families coincide:
$K_A(\{(i,j)\})$ and $K_R(\{(i,j)\})$ are the same unmarked list (the
identity piece and an empty piece), so the generated sieves are equal and
\emph{the topology cannot distinguish the two covers there}. Second, the
marked records of those same covers remain distinct: they carry different
cover kinds (\textsc{acbe} versus \textsc{rpm}) and therefore different
canonical digests. The pair corroborates the central design argument of the
marks chapter: provenance survives exactly where the topology forgets, and
an audit trail that lived only in the topology would be lost at precisely
these contexts. (Appendix~A additionally derives, from the definitions, the
full list of contexts where the unmarked families coincide; the certificate
registers the singleton instances, which are the load-bearing ones.)
\end{checkclass}

\subsection{Marked-cover integrity}
\label{ch11:sec:marks}

\begin{checkclass}
\label{ch11:chk:roundtrip}
\textbf{\texttt{marked\_cover\_roundtrip\_stable}.} Serializes marked covers
to canonical JSON, rehashes, deserializes, and verifies that the
reconstructed record and its SHA-256 digest are stable across the round
trip. This is the base integrity property of the provenance layer: a digest
that did not survive its own serialization could not chain.
\end{checkclass}

\begin{checkclass}
\label{ch11:chk:pullbackprov}
\textbf{\texttt{marked\_pullbacks\_bind\_parent\_provenance},
\texttt{marked\_pullbacks\_preserve\_topology\_family}.} Pulling a marked
cover back along an inclusion creates a new record of derivation kind
\textsc{pullback}; the first check verifies the child's provenance chain
contains its exact parent digest, so derived records are cryptographically
bound to their origins. The second verifies that pullback does not mutate
the cover family: the child's members are exactly the pullback of the
parent's members, i.e. the marked layer and the topological layer move in
lockstep under (J2). Together they corroborate the marked-cover construction
of the marks chapter and the pullback-stability geometry of
Check~\ref{ch11:chk:pullback}, now at the serialization level.
\end{checkclass}

\begin{checkclass}
\label{ch11:chk:joint}
\textbf{\texttt{joint\_refinement\_is\_four\_singletons},
\texttt{joint\_refinement\_retains\_both\_parent\_digests}.} The ACBE/RPM
common refinement at $\Omega$ is built as a derived record of kind
\textsc{common refinement}. The first check verifies its members are exactly
the four singleton covers --- the marked-layer image of the transversality
identity (Check~\ref{ch11:chk:transverse}). The second verifies the derived
record carries \emph{both} generating-cover digests as parents: the joint
refinement remembers that it came from one ACBE and one RPM cover, which is
the provenance fact the joined topology itself provably cannot retain.
Corroborates Theorem NT2-4 on the topological side and the derivation-node
schema on the audit side.
\end{checkclass}

\begin{checkclass}
\label{ch11:chk:resolve}
\textbf{\texttt{marked\_derived\_records\_resolve}.} Exercises the
fail-closed deserialization contract for derived records: resolution of a
derived record requires a parent resolver; every parent digest is rehashed,
the complete acyclic parent derivation is checked, and the child's members
and mark are recomputed and compared. The check also covers the rejection
branches: a common refinement without exactly one ACBE and one RPM
generating parent, or a pullback without its exact parent, is rejected
rather than treated as evidence, as are unknown or duplicate JSON keys and
uncontrolled review states. This corroborates no theorem; it corroborates an
\emph{invariant} --- the integrity discipline that makes the marked layer
usable as audit data at all.
\end{checkclass}

\subsection{The legacy collapse checks}
\label{ch11:sec:legacy}

\begin{checkclass}
\label{ch11:chk:legacycollapse}
\textbf{\texttt{legacy\_acbe\_identity\_forces\_maximal\_sieve},
\texttt{legacy\_rpm\_split\_epi\_forces\_identity}.} The checker builds small
executable models of the two legacy failure mechanisms and confirms each.
For the ACBE clause: a sieve containing an identity is closed under
precomposition, hence contains every arrow into its target and is maximal.
For the RPM clause: a sieve containing a split epimorphism
$\varepsilon$ with section $\eta$ contains $\varepsilon \circ \eta = \id$ by
precomposition closure, hence is maximal by the first mechanism. These are
the two elementary facts that make the legacy ACBE/RPM cover clauses
contribute only maximal sieves, collapsing the legacy topology to its
jointly-surjective clause and erasing all intended modal structure.
Corroborates the legacy-collapse proposition of the baseline chapter ---
the negative result the entire successor construction exists to repair ---
and does so inside the same certificate that verifies the repair, so that
the diagnosis and the cure are checked by one artifact.
\end{checkclass}

\begin{finding}
\label{ch11:fnd:sitesuite}
At the bound commit, the finite-site suite comprises $8$ tests and passes in
full; the certificate's aggregate flag
\texttt{all\_required\_checks\_pass} is \emph{true} over the twenty-three
registered checks above. The distinction between the $8$ tests and the $23$
checks is granularity, not content: a test is a unit-test method that may
assert many registered checks (for example, one test sweeps all pullback
verifications for both partitions), while a check is a named line of the
certificate's expected-result block. The certificate, not the test count,
is the audit surface.
\end{finding}

\subsection{Quantifier domains and counterfactual failures}
\label{ch11:sec:domains}

Exhaustiveness claims are only as strong as their stated domains. The
following table records, for each registered check group, the exact domain
enumerated and what a failure would have falsified. Domain sizes use the
site constants: $16$ contexts, $81$ inclusion arrows, $2$ partitions with
$2$ members each, $16$ representables, $2$ witnesses, $4$ singleton
contexts.

\begin{longtable}{@{}lll@{}}
\toprule
Check group & Domain exhausted & A failure would falsify \\
\midrule
\endhead
\bottomrule
\endfoot
cardinalities & the constructed $\Omega$ and $L$ &
the standing data themselves \\
proper global covers & $2$ partitions $\times$ $2$ members at $\Omega$ &
nonmaximality of the generated sieves \\
jointly exhaustive & $2$ partitions $\times$ $16$ contexts &
the coverage-basis lemma (union of pieces) \\
pullback stable & $2$ partitions $\times$ $81$ arrows $\times$ $2$
members & basis axiom (J2); Theorem NT2-1 \\
representables are sheaves & $16$ representables $\times$ $2$ topologies
$\times$ all generating covers at all $16$ contexts &
Theorem NT2-2 (subcanonicity) \\
witness functoriality & $2$ witnesses $\times$ $81$ arrows (identities
and all composable pairs) & well-definedness of the witnesses \\
four sheaf verdicts & $2$ witnesses $\times$ $2$ topologies $\times$ all
generating covers at all contexts & Theorem NT2-3 (distinctness), in
either direction \\
transverse refinement & $4$ index pairs; refinements at all $16$
contexts & Theorem NT2-4; the composition theorem's input identity \\
singleton coincidence pair & $4$ singleton contexts $\times$ both
families, unmarked and marked & the marks-necessity argument \\
marked integrity trio & serialization round trips, pullback derivations
along arrows, family preservation & the provenance-layer invariants \\
joint refinement pair & the derived record at $\Omega$ and its parent
set & the derivation-node schema; NT2-4 at the marked layer \\
derived-record resolution & resolution and all rejection branches &
the fail-closed contract \\
legacy collapse pair & both executable collapse mechanisms &
the collapse diagnosis motivating the successor \\
\end{longtable}

Reading the third column as a whole makes the certificate's structure
visible: the checks are not a grab-bag but a covering of the monograph's
finite claims --- standing data, coverage basis, subcanonicity,
distinctness, transversality, marks, and the legacy diagnosis --- such
that each finite theorem has at least one check whose failure would
falsify it in the implemented model.

\section{The modality-composition verifier}
\label{ch11:sec:composition}

The second executable layer,
\texttt{modality\_composition/verify\_modality\_composition.py}, is
independent of the finite-site checker and targets Chapter~10 specifically:
the explicit reflector formula, its unit, both flattening/regrouping maps,
and --- crucially --- the identification of the joined topology. Its
registered invocation is
\texttt{python docs/working/project\_coherence/gate1/noetic\_topos2/}
\texttt{modality\_composition/verify\_modality\_composition.py}, with a
scoped unit-test discovery command over
\texttt{tests/noetic\_topos2\_modality\_composition}.

\subsection{Exact topology closure}
\label{ch11:sec:closure}

The verifier does not sample sieves and does not trust the pen-and-paper join
argument. It computes both relevant topologies \emph{exactly}: it enumerates
every sieve on every context, seeds the covering collection with the
generated basis sieves of each presentation, and then saturates three
closure rules --- maximality (J1), pullback stability (J2), and local
character (J3) --- to a fixed point. The two presentations are computed
independently: once seeded from the ACBE and RPM bases jointly (the join
presentation), and once seeded from the singleton basis (the singleton
presentation). The run then verifies that the two fixed points are
\emph{equal} as sets of covering sieves, and checks the three Grothendieck
axioms once more on each computed result, independently, so that agreement
cannot be an artifact of a shared saturation bug reaching the same wrong
fixed point from a shared seed.

The enumeration sizes are worth deriving, because they are checkable by hand
and pin the claim that the run is exhaustive rather than sampled.

\begin{remark}[Deriving the registered counts]
\label{ch11:rem:counts}
A sieve on a context $V$ in the poset $L$ is a downward-closed set of
subcontexts of $V$, i.e. a down-set in the Boolean lattice
$\mathcal{P}(V)$. The number of down-sets of the Boolean lattice on $n$
points is the Dedekind number $M(n)$, with $M(0) = 2$, $M(1) = 3$,
$M(2) = 6$, $M(3) = 20$, $M(4) = 168$. Summing over the sixteen contexts by
size:
\[
1 \cdot M(4) + 4 \cdot M(3) + 6 \cdot M(2) + 4 \cdot M(1) + 1 \cdot M(0)
= 168 + 80 + 36 + 12 + 2 = 298
\]
sieves in all --- the registered total. For the joined (equivalently
singleton) topology, a sieve on $V$ covers iff it contains every singleton
of $V$ (Chapter~10, join corollary). Such sieves automatically contain
$\varnothing$ and are determined by a down-set of the subcontexts of size at
least two. Counting per context size: at $|V| = 0$ both sieves on
$\varnothing$ cover (the null cover generates the empty sieve, and its
supersieve is maximal), giving $2$; at $|V| = 1$ only the maximal sieve,
giving $1$ each; at $|V| = 2$ the sieve of proper subcontexts and the
maximal sieve, giving $2$ each; at $|V| = 3$ a free down-set choice over the
three $2$-subsets and the top, giving $9$ each; at $|V| = 4$ a count over
down-sets of the size-$\ge 2$ subsets of $\Omega$ gives $114$. The total is
\[
1 \cdot 2 + 4 \cdot 1 + 6 \cdot 2 + 4 \cdot 9 + 1 \cdot 114
= 2 + 4 + 12 + 36 + 114 = 168
\]
covering sieves --- the registered count, reported identically for both
presentations. The arrow count $81 = \sum_V 2^{|V|}$ was derived in
Check~\ref{ch11:chk:cardinalities}. These derivations are corroboration
arithmetic for this monograph; the registered numbers themselves come from
the verifier's enumeration. Appendix~A tabulates the per-context values.
\end{remark}

\subsection{What the composition run checks}
\label{ch11:sec:comprun}

The deterministic run verifies, in order:

\begin{enumerate}
\item \textbf{Base combinatorics.} All $4$ atoms, $16$ contexts, and $81$
inclusion arrows are constructed and counted.
\item \textbf{The sieve universe.} All $298$ sieves across the sixteen
contexts are enumerated.
\item \textbf{Topology equality.} The join presentation and the singleton
presentation saturate to the same $168$ covering sieves --- the executable
content of the join lemma $J_A \vee J_R = J_{\mathrm{atom}}$.
\item \textbf{Axioms on both results.} The Grothendieck axioms (J1)--(J3)
are verified independently on each of the two computed topologies.
\item \textbf{Composite decompositions.} At every context, both index-level
decompositions of the nested product --- the ACBE-then-RPM expansion and the
RPM-then-ACBE expansion --- are computed and compared with the atomic
product, corroborating the objectwise computation in the composition
theorem's proof.
\item \textbf{Natural isomorphisms on an adversarial fixture.} Both
canonical natural isomorphisms are checked on a constant two-section
presheaf whose empty-context value is \emph{deliberately nonterminal}. This
fixture exists to exercise the $(R_P F)(\varnothing) = \mathbf{1}$ clause:
a wrong implementation that carried $F(\varnothing)$ through the reflector
would satisfy most checks on well-behaved presheaves and fail exactly here,
at the null cover and at every empty overlap.
\item \textbf{Structured fixtures.} A bit-assignment presheaf (the witness
shape) and all $16$ representables are run through the same comparisons.
\item \textbf{The full battery.} Inverse laws for the flattening and
regrouping maps, restriction naturality (naturality in $V$), commutation
naturality (naturality in $F$ of the composed comparison), and unit
compatibility are verified across $18$ executable fixtures and $3{,}012$
restriction squares.
\end{enumerate}

Each item corroborates a specific labeled step of Chapter~10: items (3) and
(4) the join lemma; item (5) the objectwise flattening computation; items
(6)--(8) the naturality, inverse, and unit-compatibility paragraphs of the
composition theorem and the well-definedness proposition of the reflector.

\subsection{The fixture battery}
\label{ch11:sec:fixtures}

The $18$ executable fixtures decompose as $16 + 1 + 1$: the sixteen
representable presheaves, the constant two-section presheaf with
nonterminal empty-context value, and the bit-assignment presheaf. The
three families are chosen to stress different failure surfaces, and it is
worth recording why each earns its place.

\emph{The representables} are the fixtures on which everything should be
\emph{inert}: they are sheaves for both factor topologies
(subcanonicity), so the units and both composites should act by
isomorphisms, and any fixture-dependent bug that breaks canonicity ---
an indexing error that permutes coordinates, say --- shows up as a
non-isomorphism or a failed naturality square on some $yW$. Sixteen of
them cover every shape of context, including the empty representable
$y\varnothing$ (a singleton exactly at $\varnothing$), which probes the
null-cover conventions from the sheaf side.

\emph{The constant two-section presheaf} is the designed adversary for
the empty-product clause, as discussed in Chapter~10: its empty-context
value has two elements, so any implementation that carries
$F(\varnothing)$ through the reflector, or that treats empty overlaps as
imposing real constraints, fails on it and on nothing else so cheaply.
It also exercises value \emph{shrinkage} at $\varnothing$ simultaneously
with value \emph{growth} at $\Omega$ (the reflector value there is a
product of two-element sets), so restriction squares crossing the empty
context are genuinely two-sided.

\emph{The bit-assignment presheaf} has the witnesses' shape: values of
different cardinalities at different contexts, coordinate-forgetting
restrictions, and sensitivity to which blocks are complete. On this
fixture the composite decompositions are nondegenerate --- nested and
flat products differ as concrete data structures --- so the flattening
and regrouping maps are tested where they actually move something.

Across the battery, the run verifies the inverse laws
($\psi \circ \varphi = \id$ and $\varphi \circ \psi = \id$ objectwise),
restriction naturality, commutation naturality, and unit compatibility,
totaling $3{,}012$ restriction squares --- the registered count of
commuting-square instances discharged; the per-fixture breakdown is an
implementation detail of the verifier and the monograph cites only the
registered total.

\subsection{Rerunning the layer}
\label{ch11:sec:rerun}

The composition layer's registered commands are reproduced in the
artifact and in Chapter~10's source document; the deterministic verifier
is invoked as a single script with no arguments, and the scoped suite by
standard unit-test discovery over
\texttt{tests/noetic\_topos2\_modality\_composition}. A reviewer's
workflow is: check out the bound commit, run both commands, and diff the
verifier's report against the registered expectations. Any discrepancy
is, by the fail-closed discipline, a finding about the artifact set ---
either a regression or an environment sensitivity --- and not something
to be waved through.

\begin{finding}
\label{ch11:fnd:compsuite}
At the bound commit, the modality-composition suite comprises $6$ tests and
passes in full, and the deterministic verifier run completes with every
check above reporting success.
\end{finding}

\section{The Lean layer: scope and explicit non-scope}
\label{ch11:sec:lean}

The bounded proof-assistant layer is a pinned Lean Gate~1 verifier that
passes at the bound commit. Its scope is deliberately narrow, and the
narrowness is registered rather than incidental.

\begin{checkclass}
\label{ch11:chk:lean}
\textbf{Formalized (in scope).} Four-atom exhaustiveness (the atoms
enumerate $\Omega$); properness of the global cover pieces; the transverse
refinement identity $A_i \cap R_j = \{(i, j)\}$; and the witness-map lemmas
(the elementary combinatorial facts about complete-block containment that
feed the witness constructions).
\end{checkclass}

\begin{boundary}
\label{ch11:bnd:leannonscope}
\textbf{Not formalized (explicit non-scope).} The Lean layer does
\emph{not} formalize: the sixteen presheaves; any sheaf condition; topology
generation or the saturation of (J1)--(J3); the subcanonicity theorem; the
witness sheaf/non-sheaf verdicts; the reflector construction or its
universal property; the join lemma; the composition theorem; the
product-site comparison; or the legacy round trip. Any sentence of the form
``the Noetic Topos 2.0 results are formally verified'' is therefore false.
The accurate sentence is: a small set of finite combinatorial lemmas at the
base of the construction is machine-checked in Lean; the sheaf-theoretic
theory above them is informally proved and exhaustively executed, and its
proof-assistant formalization is a registered open obligation
(Chapter~12).
\end{boundary}

\section{The certificate's own claim boundary}
\label{ch11:sec:boundary}

The certificate closes with a claim boundary of its own, and it governs
every number in this chapter. Quoting the registered artifact
\texttt{FINITE\_MODEL\_CERTIFICATE\_v0.1.json} verbatim:

\begin{quote}
``This is exhaustive execution over one sixteen-object finite context site.
It is not proof-assistant verification, external peer review, empirical
validation, or a product-claim license.''
\end{quote}

\begin{boundary}
\label{ch11:bnd:governing}
Unpacking the four negations, each against the temptation it blocks:
\begin{enumerate}
\item \emph{Not proof-assistant verification.} The Python layers are
ordinary programs; their correctness is not machine-checked, and the Lean
scope above excludes the sheaf theory. Exhaustiveness over the model does
not transfer trust to the model's formalization.
\item \emph{Not external peer review.} No independent category theorist has
reviewed the proofs or the code. The composition result in particular
remains an \textsc{internal proof candidate / review required} artifact.
\item \emph{Not empirical validation.} Nothing here measures anything. The
ACBE and RPM coordinates are formal marks; no execution result bears on
whether either coordinate detects any cognitive, causal, safety, or
alignment property of any system.
\item \emph{Not a product-claim license.} The certificate's passing state
may not be cited, in any register of the program, as grounds for a
capability, quality, compliance, or assessment claim about any product. The
artifact set remains Ring~4 and \textsc{not licensed}.
\end{enumerate}
What the execution \emph{does} establish is narrower and still valuable:
that the registered finite model, as implemented, satisfies every one of the
enumerated properties without exception, over its entire state space; that
the prose theorems and the executable model agree everywhere both speak;
and that the negative expectations (the two \emph{false} sheaf verdicts,
the rejection branches of the marked-record resolver, the two legacy
collapse mechanisms) are checked as first-class results rather than assumed.
\end{boundary}
